
\section{Related Work}
% ──────────────────────────────────────────────────────────────────────────────

Efforts to taxonomise XAI methods have produced a substantial but fragmented
body of work. Arrieta et al.~\cite{Arrieta2020XAI} established the foundational
vocabulary --- distinguishing interpretable-by-design models from post-hoc
explanation methods, and scoping the core challenges toward responsible AI ---
that most subsequent taxonomies inherit. Schwalbe and
Finzel~\cite{schwalbe2024comprehensive} extended this by synthesising more than
fifty XAI surveys into a unified framework organised around problem definition,
explanator properties, and evaluation metrics, finding that high-level properties
such as faithfulness and stability recur broadly but their operationalisations
diverge sharply by task and model type. Speith~\cite{speith2022review} identified
a structural reason for this divergence: taxonomies organised around how a method
works versus what it produces classify the same methods differently, making
principled cross-study comparison unreliable. Taken together, these works
establish that the field lacks not just shared metrics but a shared vocabulary
for organising evaluation criteria --- a problem that is especially acute once
XAI moves into high-stakes clinical settings.

Taxonomies built specifically for medical AI have begun to address this, but
remain incomplete. Jin et al.~\cite{jin2023guidelines} proposed clinical
evaluation guidelines that explicitly separate technical faithfulness from
clinical usefulness, offering the clearest domain-specific taxonomy to date.
Holzinger et al.~\cite{holzinger2022information} argued that
robustness under distributional shift is a persistent gap that existing medical
XAI taxonomies handle inconsistently. Tjoa and
Guan~\cite{Tjoa2021MedicalXAI} surveyed the broader landscape of medical XAI
methods, but without a structured evaluation taxonomy, leaving the question of
how to assess explanation quality largely open. More recently,
Varghese et al.~\cite{Varghese2025OnDE} applied Borda count and multiple
correlation techniques to rank competing XAI metrics for image classification,
providing empirical grounding for treating faithfulness and sensitivity as
primary comparison axes --- though their analysis does not extend to
cross-domain alignment within medicine.

The metrics literature exposes the consequences of this taxonomic gap.
Rosenfeld~\cite{rosenfeld2021better} showed that most evaluation criteria focus
too narrowly on algorithmic fidelity, systematically omitting human-centric
concerns. Kadir et al.~\cite{kadir2023evaluation} confirmed this through a
PRISMA-based review: sensitivity, faithfulness, and robustness dominate the
quantitative literature, while human-centred assessment lags behind.
Without a taxonomy that structures which metrics belong
at which layer of evaluation, Phuong et al.~\cite{phuong2023benchmarking} found
that XAI toolkits produce non-comparable results, and Klein et
al.~\cite{Klein2024NavigatingTM} showed that metric selection remains
effectively arbitrary even in state-of-the-art benchmarks.
Mirzaei et al.~\cite{mirzaei2023xai} proposed a top-down
selection framework as a partial remedy, but without grounding in a medical
taxonomy, the framework does not transfer cleanly to clinical contexts.

What is missing across all of this work is a taxonomy constructed specifically
to organise XAI evaluation criteria across medical sub-domains, validated
against real usage patterns rather than derived deductively from general XAI
literature. No existing framework jointly covers oncology, cardiology,
neurology, and radiology, accounts for the three-layer structure of
mathematical faithfulness, system robustness, and human-centred trust, or
analyses metric co-occurrence and redundancy at a cross-domain level. The
present study fills this gap directly.
